\section{Introduction}
\label{sec:intro}

Predicting several future tokens at once (multi-token prediction, MTP) has
become a common ingredient in large-scale language modeling, used to improve
sample efficiency and to enable self-speculative decoding
\citep{gloeckle2024mtp,deepseekv3}. The auxiliary objective is double-edged: a
growing body of work reports that adding it can \emph{degrade} the primary
next-token objective, particularly at smaller scales and longer horizons
\citep{gloeckle2024mtp,zuhri2025top}. Asked why, practitioners reach for the
intuition that pervades multi-task learning: the two losses must pull the shared
parameters in \emph{conflicting directions}, and the remedy is to detect and
neutralize that conflict, typically from the sign of a gradient cosine
\citep{yu2020pcgrad,wang2021gradvaccine}. This intuition has been questioned in
auxiliary-task learning: \citet{jiang2023forkmerge} report that gradient-conflict
coordination can overlook auxiliary-target generalization, so the cosine signal
does not by itself determine negative transfer. We ask what that signal actually
measures at a language model's shared output head.

This paper is a measurement study of that intuition for the case of an auxiliary
multi-token objective sharing a model's tied output projection. Our central
finding is diagnostic: \textbf{the cosine conflict signal is misleading here.} We
instrument the gradient each loss deposits on the output projection and examine
it \textbf{row by row} (one row per vocabulary item). Two facts emerge, both from
data already in hand. First, the per-row gradients are not opposed; they are \emph{aligned}. We measure this on the rows that receive a target in the measurement batch, the ${\approx}8\%$ that are not scalar multiples by construction (Section~\ref{sec:aligned}). On those rows the median per-row cosine is ${\approx}0.53$, and ${\approx}3$--$4\%$ are opposed. Over all rows the median is $+1$ and $98.4\%$ have $|\cos| > 0.3$. On the matched active-row denominator, $83\%$ of active rows have $\cos > 0.3$; a control that measures CE against an unrelated L1 loss on the same rows reads $1.0\%$. Either way the per-row directions are
predominantly aligned, in sharp contrast to the near-zero aggregate. Second,
the \emph{aggregate} cosine over the flattened matrix is nonetheless near zero,
and tracking it over training shows the two metrics \textbf{start together} (both
$\approx +0.99$ at initialization) and then \textbf{decouple}: within ${\sim}200$
steps the aggregate falls below zero while the per-row median stays at
$\approx +1$. By count, almost no rows are opposed. But the aggregate cosine is magnitude-weighted, so we instrumented the per-row gradient norms directly to see where the weight lies. The two losses load their magnitude on the \emph{same} rows: the norm-profile cosine, meaning the cosine of the two per-row norm vectors, is $0.98$. Yet a tiny high-norm opposed minority, ${\approx}0.3\%$ of rows, carries $60$--$69\%$ of the mass across optimizers. It is this minority that pulls the flattened sum negative. The
alignment is real and holds for the vast majority of rows; the aggregate is a
\emph{norm-weighted cancellation} the cosine cannot localize.

We frame this norm-weighted opposition as the structure a cosine diagnostic
cannot see (Figure~\ref{fig:schematic}), and we relate it to the observed
degradation without claiming to have isolated its cause. Adding the shared MTP objective degrades next-token loss by
${\approx}0.39$ nats (${\approx}10\%$) at 124M, a large and unambiguous effect
(Welch's unequal-variance $t$-test, $p\approx 8\times10^{-5}$), yet it co-occurs with aligned, not
conflicting, per-row gradients (the degradation is measured on the full
30{,}517-step run; per-row alignment on the diagnostic runs above). Three further
observations, all from committed
runs, are consistent with a norm-weighting/capacity account rather than a
directional-conflict one. (i) Conflict-targeted gradient surgery does not recover
the baseline even though the opposition it targets is real and mass-dominant: the
opposed rows are negligible by count ($0.3\%$) but carry the majority of the
gradient-magnitude mass ($60$--$69\%$), and per-row Gram--Schmidt on the head
(Appendix~\ref{app:surgery}) removes exactly that opposition, yet the degradation
does not move. The opposition the aggregate cosine detects is mass-dominant and causally inert (single-seed; Appendix~\ref{app:surgery}). (ii) Detaching the auxiliary's gradient from
the shared head makes the damage \emph{worse}, not better. The probe forms the
auxiliary logits as $\mathrm{sg}(W)h$, so the auxiliary still trains the trunk and
only the head's auxiliary gradient is removed; it is confounded because the head can
no longer co-adapt to the representations the auxiliary shapes, and we treat it as
suggestive only. (iii)
An auxiliary that avoids the shared head (a layer-split contrastive auxiliary)
reduces but does \emph{not} eliminate the gap. These auxiliaries change the objective
as well as its capacity, so the capacity ingredient is not isolated; the untested comparison that would isolate it is separate-head MTP (Section~\ref{sec:limitations}). The figures above are for the contrastive auxiliary: its standard variant is demonstrably short of the CE-only baseline ($\Delta{+}0.039$, $p \approx 0.010$), while a hyperparameter-tuned variant sits below our detection floor, the smallest difference these seed counts can resolve ($\Delta{+}0.018$, $p \approx 0.19$), and is inconclusive. We rest the capacity conclusion on the standard variant.

An exact rewrite of the shared-head objective ties these three observations
together. The objective is cross-entropy against a next-and-future-token mixture label, in effect label smoothing toward the future (an analogy: conventional label smoothing mixes toward the uniform distribution, whereas this mixes toward future-token targets). Its minimizer is therefore not the next-token distribution, which makes the degradation a plausible cost of the target rather than of the optimization path. On that reading
the norm-weighting/capacity observations are subsumed rather than displaced: surgery
would be inert because reprojecting gradients does not move an optimum (i; exact for symmetric projection, empirical for our head-only variants, Section~\ref{sec:discussion}), and
separating capacity would help because it stops the two objectives from sharing one
label (iii). A zero-parameter estimate of the implied KL cost accounts, raw,
for $62$--$91\%$ of the matched within-sweep gap ($0.3715$) across estimator
variants; the coverage-extrapolated full-mixture variant lands at ${\approx}88$--$89\%$, leaving a residual we name
rather than assign.

These findings carry a methodological point beyond MTP. Cosine-based conflict
detection is standard practice in multi-task and auxiliary-learning optimization
\citep{yu2020pcgrad,wang2021gradvaccine,du2018auxsim}, and it is the wrong
instrument in this regime, misleading in two opposite directions depending on the
optimizer. Under Muon the head aggregate sits near zero (${-}0.08$ to ${-}0.19$
across measurement batches), which a practitioner treating near-zero as benign reads
as ``no interference'' while the loss degrades regardless. Under AdamW it goes
negative enough to trip a conflict flag, yet surgery targeting exactly that
opposition recovers nothing. Both readings are of the \emph{head} aggregate; a
practitioner applying a strict sign rule, or aggregating over all parameters, would
flag conflict under both optimizers (Section~\ref{sec:general}). The
decomposition that exposes the problem separates \textbf{direction} (cosine) from
\textbf{support} (where each gradient's magnitude lives). Two additional
measurements sharpen the picture and guard against over-reading it. The high
per-row alignment is \textbf{general} across the network's weight matrices (attention and MLP matrices show it too, many with strongly \emph{negative}
aggregate cosine), so the output head is not special because the metric singles
it out; it is consequential because it directly produces the logits. And the
alignment \textbf{persists when the head is untied}, so it is a property of the
output-projection gradient, not an artifact of weight tying. A full-vocabulary
reading suggests rarer tokens align more, but we show this is a construction
artifact: rare tokens are almost never active, so their rows are dominated by
the cosine-$+1$ parallel majority, and on active rows the effect vanishes. Finally,
our account is a direction-space, multi-loss
complement to \citet{godey2026lmhead}, who show that backpropagating a
\emph{single} loss through the low-rank output head destroys most of its gradient
\emph{norm}; we ask what a \emph{cosine} diagnostic reports when \emph{two} losses
share that head, and answer that their per-row directions align while a high-norm
opposed minority sets the aggregate: measured directly via the per-row gradient
norms (Section~\ref{sec:mechanism}), not inferred.

We are explicit about scope. Every result is for a \textbf{single shared tied
output head}, at 124M ($n{=}5$) and 350M (single seed), undertrained, small-batch.
We do \emph{not} claim that multi-token prediction \emph{as practiced} (separate per-horizon
heads or modules, typically feeding a shared unembedding) degrades next-token modeling; establishing that would require a
separate-head experiment we identify as the primary future step (Section~\ref{sec:limitations}), along with $n{=}3$ at 350M and matched-seed
surgery. What we do establish, on data in hand, is that the standard cosine
conflict diagnostic misreads this regime, and why.

\paragraph{Contributions.} All on a controlled, reproducible GPT-2 testbed (124M
primary, 350M illustrative; FineWeb-Edu) with open code and a committed statistics
script:
\begin{itemize}
\item \textbf{A diagnostic result (Figures~\ref{fig:perrow},~\ref{fig:emergent}).}
  For an auxiliary multi-token objective sharing the tied output projection, the
  per-row gradients of the two losses are \emph{aligned in direction} on the rows
  that carry a target (active-row median cosine ${\approx}0.53$, ${\approx}3$--$4\%$
  opposed; over the full vocabulary the median is $+1$ because ${\approx}92\%$
  of rows are scalar multiples by construction), while the aggregate cosine is
  $\approx 0$, and this decoupling \emph{emerges during training} from an aligned
  start. Standard cosine-based conflict diagnostics therefore \emph{misdiagnose}
  this regime.
\item \textbf{A norm decomposition of the collapse, and a direct measurement of
  its cause (Figures~\ref{fig:normsupport},~\ref{fig:normdecomp}).} Via an exact
  identity, the aggregate would read $+0.94$/$+0.96$ if gradient magnitude were
  flat across rows but reads $-0.08$/$-0.28$: the collapse is set by per-row
  \emph{norm} structure, not by the typical per-row direction. Instrumenting the per-row norms directly (\texttt{measure\_norm\_support.py}; $n{=}3$ seeds per optimizer) settles which structure. The two losses load the \emph{same} rows: norm-profile cosine $0.98$. A high-norm opposed minority carries the majority of the mass: opposed-norm fraction $0.60$ under Muon and $0.69$ under AdamW. The
  aggregate is a norm-weighted cancellation, not disjoint support.
\item \textbf{A calibrated, general characterization
  (Figures~\ref{fig:perlayer},~\ref{fig:tokenfreq}).} The alignment is
  \textbf{general across the network} and \textbf{survives untying} the head; the
  apparent rarer-tokens-align-more effect is a construction artifact that vanishes
  on active rows. The next-token degradation is large and
  significant ($+0.39$ nats, $p\approx 8\times10^{-5}$) yet co-occurs with aligned
  gradients.
\item \textbf{Intervention evidence against a directional-conflict account
  (Figure~\ref{fig:interventions}, Appendix~\ref{app:surgery}).}
  Conflict-targeted surgery does not recover the baseline; stop-gradient makes the
  damage worse; and the standard capacity-separated auxiliary reduces but does not eliminate
  the gap (the standard variant demonstrably so; the tuned variant inconclusive), each reported with
  its statistical test and caveats.
\item \textbf{An account of what the shared head actually optimizes, and a
  zero-parameter test of it (Figure~\ref{fig:mixturesweep}).} An exact rewrite shows
  the shared-head objective is $1.75\times$ cross-entropy against a next-and-future-token
  mixture label, so its optimum moves and the next-token cost is bounded by
  $\mathrm{KL}(P_1 \Vert q^{*}) \leq \log 1.75$. Estimating that KL from the CE-only
  model's own distributions, with no fitted parameters, reproduces the measured
  degradation across an anchored auxiliary-weight sweep, matching within
  $0.1\sigma$ at the lower weights and falling $1.4$--$1.8\sigma$ short at full
  weight. This makes the shifted optimum, rather than gradient conflict, the leading
  account of the loss increase, and it explains why conflict-targeted surgery cannot
  help: symmetric projection leaves a shifted optimum stationary, and the head-only
  variants we test are empirically inert. We name the residual
  rather than claim the account is confirmed.
\item \textbf{A bridge to the output-layer bottleneck literature.} We position the
  result as the direction-space, multi-loss complement to
  \citet{godey2026lmhead}'s single-loss, norm-space account of the output head as
  a gradient bottleneck.
\end{itemize}

\paragraph{Roadmap.} Section~\ref{sec:aligned} establishes the diagnostic
failure: on the rows that carry supervision the two gradients agree, while the
aggregate the standard instrument reports does not.
Sections~\ref{sec:emergent} and~\ref{sec:mechanism} give the mechanism, an exact
two-factor decomposition showing that a high-norm opposed minority, not disjoint
support, sets the aggregate.
Sections~\ref{sec:softlabel} through~\ref{sec:interventions} give the account of
the loss increase itself: the shared head makes the objective a cross-entropy
against a blended target, so the optimum moves, which is why gradient surgery
leaves the degradation untouched. Two further subsections carry load this summary
would otherwise hide. Section~\ref{sec:general} shows the signature is not specific
to this loss pair. Section~\ref{sec:rare} withdraws a rare-token reading we can show
is a construction artifact, and reports the 350M scale check.


\begin{figure}[t]
\centering
\includegraphics[width=\linewidth]{figures/fig0_schematic.pdf}
\caption{\textbf{Most rows agree; a high-norm minority sets the sign} (schematic; illustrative proportions, not plotted data). \emph{Left:} row by row (one
row per vocabulary token), the CE and MTP gradients point almost the same way for
$99.7\%$ of rows (per-row cosine $\approx +1$); a tiny high-norm minority (red) is
opposed. \emph{Middle:} both losses deposit their gradient \emph{magnitude} on the
\emph{same} rows (norm-profile cosine $\approx 0.98$), and one high-norm row (red)
carries most of the mass. \emph{Right:} the typical per-row cosine reports
alignment, while the magnitude-weighted aggregate, the standard conflict
diagnostic, is dragged negative by the opposed high-norm minority and is
misinterpreted as conflict. The aggregate is not wrong about the update, but it
cannot say \emph{where} the interference lives: it is set by per-row norms, not by
the directions it averages. The $99.7\%$ on the left is the fraction of rows that are
\emph{not} opposed, and the $\approx +1$ is dominated by the ${\approx}92\%$ of
rows that receive no target in the batch and are exact scalar multiples. On the
${\approx}8\%$ that do carry supervision the median cosine is ${\approx}0.53$
with ${\approx}3$--$4\%$ opposed (Section~\ref{sec:aligned}); those are the quantities the
paper's claims rest on.}
\label{fig:schematic}
\end{figure}

% ---------------------------------------------------------------------------
